% !TEX root = main.tex
\section{A Dynamical View on Power Difference Problem}~\label{sec:a-dynamical-view-on-power-difference-problem}
In \cref{theorem:inductive-formula}, we can find the exponent of $\alpha$ is always $n$ wihle exponent of $\beta$ is variable with reduce sequence. If we make the exponent of $\alpha$ is also a accumulation of a reduce sequence, then we have
\begin{theorem}~\label{thm:power-diff-system-exists}
	Given two arguments $\alpha, \beta \in \mathbb{N}_{\ge2}$ such that $\gcd(\alpha, \beta) = 1$.\newline
	Given two sequence $\{a_i\}_{i=1}^{n}, \{b_i\}_{i=1}^{n}\subset\mathbb{N}$ and $a\in \mathbb{N}$, then $\exists \{\gamma_i\}_{i=1}^{n}\subset\mathbb{N}$ and $b \in\mathbb{N}$, such that
	\[ b\cdot\beta^{\mathcal{B}_n} - a\cdot\alpha^{\mathcal{A}_n} = \sum_{k=0}^{n-1}\alpha^{\mathcal{A}_k}\beta^{\mathcal{B}_{n-1}^{k}}\gamma_k. \]
\end{theorem}
\begin{proof}
	Similar with \cref{theorem:inductive-formula}. Let $\gamma_k = a_k\cdot\beta^{\mathcal{B}_k^{k-1}} - a\cdot\alpha^{\mathcal{A}_{k}^{k-1}}$, then we have 
	\[ a_k\cdot\beta^{\mathcal{B}_k^{k-1}} = a\cdot\alpha^{\mathcal{A}_{k}^{k-1}} + \gamma_k. \]
	Thus, just let $b = a_n$, the theorem holds.
\end{proof}
\begin{definition}~\label{def:derivation-of-power-diff}
	We denote $ \mathfrak{U}_{n} = \sum_{k=0}^{n-1}\alpha^{\mathcal{A}_k}\beta^{\mathcal{B}_{n-1}^{k}}\gamma_k$ in \cref{thm:power-diff-system-exists}.
	We have $\mathfrak{U}\in\partial_{*}(\alpha^{\mathcal{A}_{*}}, \beta^{\mathcal{B}_{*}})$ follows conclusions in \cref{sec:a-kind-of-derivation-on-groupoid} and its kernel is exactly $\gamma_{k}$.
\end{definition}
\begin{remark}
	Notice that the system in \cref{thm:power-diff-system-exists} is built from reduce sequence, thus we have no result like $\beta \nmid b$ and $\beta\nmid \mathfrak{U}_n$.
\end{remark}
\begin{corollary}~\label{cor:power-diff-to-dynamical-system}
	Suppose $\alpha^{A} - \beta^{B} = C\in\mathbb{N}$. Then $\exists t\in\mathbb{N}$ and sequence $\{a_i\}_{i=1}^{t} \subset\mathbb{N}, \{b_i\}_{i=1}^{t} \subset\mathbb{N}$, $\gamma \in\mathbb{N}_{\ge1}$, such that $A = \mathcal{A}_t$, $B = \mathcal{B}_t$, and $C\mid \mathfrak{U}_{t}$.
\end{corollary}
\begin{proof}
	Extend sequences $\{a_i\}$ and $\{b_i\}$ as $t$-periodic with $a_i = a_{i\fmod t}$, $b_i = b_{i\fmod t}$, and let $\gamma' = 1$. Suppose the derivation is $\mathfrak{V}$, then $\mathfrak{V}\in\partial(\alpha^{\mathcal{A}_{*}}, \beta^{\mathcal{B}_{*}})$.
	From \cref{thm:cycles-derive-cycle}, and alternative with \cref{theorem:any-cycle}, we have $b = a$ and
	\[ a = \frac{\mathfrak{V}_{t}}{\beta^{\mathcal{B}_t} -  \alpha^{\mathcal{A}_t}} = \frac{\mathfrak{V}_{t}}{C} \in\mathbb{Q}. \]
	Suppose $d = \gcd(C, \mathfrak{V}_{t})$, then let $\gamma = \frac{C}{d}\cdot\gamma'$, then we have $\mathfrak{U} = \frac{C}{d}\cdot\mathfrak{V}$ such that $a = \frac{\mathfrak{U}_{t}}{C} \in \mathbb{N}$. Thus $C\mid \mathfrak{U}_{t}$.
\end{proof}
\cref{cor:power-diff-to-dynamical-system} provide a view that given $C$, we can find the solution $A$ and $B$ by search the partition of exponents.
